% ==============================================================================
% lab3_incident_simulation.tex — UE SIS589
% Travaux Pratiques 3 : Simulation Tabletop — Ransomware LiZbiZ-SIS
% ==============================================================================

\chapter{TP~3~— Simulation Tabletop~: Ransomware sur LiZbiZ-SIS}
\label{lab:tabletop}

\epigraph{\itshape « On ne s'improvise pas gestionnaire d'incident.
  On s'y prépare. »}
         {--- NIST SP~800-61}

% ── Présentation ──────────────────────────────────────────────────────────────
\section*{Présentation du TP}

\begin{table}[H]
\centering\small
\rowcolors{2}{TableauPair}{TableauImpair}
\begin{tabular}{p{3.5cm}p{10cm}}
\toprule
\textbf{Durée} & 4~heures (2~h simulation + 2~h rédaction) \\
\textbf{Format} & Exercice Tabletop en équipes de 4~à~5 étudiants \\
\textbf{Objectifs} & Piloter une réponse à incident ransomware en temps
  contraint~; appliquer le cycle ISO~27035~; rédiger un rapport exécutif
  et un plan de remédiation~; conduire le RETEX \\
\textbf{Pré-requis} & Chapitres~8, 9 et~10~; Lab~1 et Lab~2 complétés \\
\textbf{Rôles} & Chaque équipe désigne~: 1~Incident Manager,
  1~Analyste Technique, 1~Responsable Communication,
  1~Coordinateur Légal, 1~Secrétaire (log d'incident) \\
\textbf{Livrable} & Rapport d'incident (15~pages)~; présentation
  orale de 10~minutes \\
\bottomrule
\end{tabular}
\end{table}


% ==============================================================================
\section{Scénario initial~: l'alerte de 07h13}
\label{lab3:scenario-initial}
% ==============================================================================

\begin{boxalert}
\textbf{Contexte~— Lundi 15~mai~2026, 07h13~UTC}

L'analyste N1 de permanence reçoit simultanément~:

\begin{itemize}
  \item Une alerte SIEM critique~: \texttt{RULE 100201 — Ransomware
        Indicator — Mass file rename (.enc extension)} depuis
        srv-erp-01 (10.0.2.10, Windows Server~2022).
  \item Un appel téléphonique du responsable comptabilité~:
        \og{}Tous mes fichiers ont une extension bizarre (.enc),
        je ne peux plus les ouvrir, et il y a un fichier
        READ\_ME\_DECRYPT.txt sur mon bureau.\fg{}
  \item Un email automatique du système de sauvegarde~:
        \og{}Erreur~: impossible d'accéder au dossier
        \texttt{\textbackslash\textbackslash srv-backup-01\textbackslash backups}.
        Accès refusé.\fg{}
\end{itemize}

\textbf{Informations disponibles à T+0~:}
\begin{itemize}
  \item Il est 07h13. La DSI ouvre à 08h00. Le RSSI est en déplacement
        à Douala. Le DG n'est pas encore arrivé.
  \item Le SIEM affiche 847~alertes en 12~minutes sur srv-erp-01.
  \item L'IP 185.220.101.5 (sortante) apparaît dans les logs
        firewall depuis 03h47.
  \item Aucun playbook ransomware n'existe.
  \item Les sauvegardes de la semaine précédente se trouvent sur
        srv-backup-01 (Windows Server~2019), sur le même VLAN~20.
\end{itemize}
\end{boxalert}


% ==============================================================================
\section{Phase~1 (T+0 à T+30~min)~: Détection et triage}
\label{lab3:phase1}
% ==============================================================================

\textbf{Durée allouée~: 30~minutes de simulation.}

\begin{boxexo}
\textbf{Travail Phase~1 — Questions à traiter en équipe~:}

\begin{enumerate}
  \item \textbf{Qualification (Incident Manager + Analyste)~:}
        Classifiez cet incident selon la grille du chapitre~8
        (tableau~\ref{tab:gravite-incidents}). Justifiez en 3~phrases.
        Qui activez-vous en premier~?

  \item \textbf{Premières décisions (Incident Manager)~:}
        Listez les 5~premières actions à mener dans les 30~minutes,
        dans l'ordre de priorité. Pour chacune~: qui la mène~?
        durée estimée~? risque si non faite~?

  \item \textbf{Isolation (Analyste Technique)~:}
        Sur quels systèmes agissez-vous en premier~? Décrivez
        techniquement comment vous isolez srv-erp-01 du réseau
        sans l'éteindre. Pourquoi ne pas l'éteindre immédiatement~?

  \item \textbf{Communication immédiate (Responsable Communication)~:}
        Rédigez le message de 3~phrases envoyé à TOUS les employés
        de LiZbiZ-SIS dans les 10~premières minutes. Objectif~:
        prévenir la propagation sans créer de panique.

  \item \textbf{Aspect légal (Coordinateur Légal)~:}
        La base de données clients (50~000 dossiers incluant noms,
        emails, historiques de commandes et transactions mobile money)
        est sur srv-erp-01. Quelles obligations légales s'appliquent~?
        Quel délai~? Quel destinataire~?
\end{enumerate}
\end{boxexo}

\begin{boxlizbiz}
\textbf{Inject~1 — T+15~min}

L'analyste découvre dans les logs firewall~:
\begin{itemize}
  \item Connexions SMB en masse depuis srv-erp-01 vers srv-backup-01
    entre 03h47 et 04h22 (volume~: 2,3~Go transféré).
  \item Connexion RDP réussie depuis 196.207.12.88 vers srv-erp-01
    le 14~mai à 22h41 (la veille), user~: \texttt{administrator}.
  \item Connexion sortante HTTPS vers 185.220.101.5~:443 établie
    depuis srv-erp-01, durée~: 7h34~min.
\end{itemize}
\end{boxlizbiz}

\begin{boxexo}
\textbf{Questions Post-Inject~1~:}

\begin{enumerate}
  \item Révisez votre évaluation initiale. L'incident est-il plus
        grave que prévu~? Pourquoi~?
  \item La connexion RDP du 14~mai à 22h41 change-t-elle votre
        analyse~? Quel est le vecteur d'accès initial probable~?
        Quel TTP MITRE~?
  \item Srv-backup-01 est-il compromis~? Sur quelle base~?
        Quelle décision prenez-vous concernant les sauvegardes~?
\end{enumerate}
\end{boxexo}


% ==============================================================================
\section{Phase~2 (T+30 à T+2h)~: Investigation et confinement}
\label{lab3:phase2}
% ==============================================================================

\textbf{Durée allouée~: 45~minutes de simulation.}

\begin{boxlizbiz}
\textbf{Inject~2 — T+35~min}

L'analyste a pu se connecter localement à srv-erp-01 (pas SSH, pas RDP —
connexion console directe). Il observe~:
\begin{itemize}
  \item Le processus \texttt{svchosts.exe} (noter le 's' final) tourne
    dans \texttt{C:\textbackslash ProgramData\textbackslash Microsoft\textbackslash
    svchost\textbackslash}.
  \item Une tâche planifiée nommée \texttt{WindowsUpdater} exécute
    \texttt{C:\textbackslash ProgramData\textbackslash upd.vbs} toutes les
    5~minutes.
  \item Le volume des données chiffrées~: 47~Go sur les partages réseau
    cartographiés (\texttt{Z:\textbackslash} = srv-backup-01).
  \item Le fichier \texttt{READ\_ME\_DECRYPT.txt}~:
    \og{}\textit{Your files are encrypted with AES-256. Send 5~BTC to
    [adresse] within 72h or data will be published online.}\fg{}
  \item Connexion active (C2) toujours établie vers 185.220.101.5.
\end{itemize}
\end{boxlizbiz}

\begin{boxexo}
\textbf{Travail Phase~2~:}

\begin{enumerate}
  \item \textbf{Confinement (Analyste Technique)~:}
    Listez les 6~actions techniques de confinement dans l'ordre.
    Pour chacune~: commande exacte (PowerShell / CLI switch) ou
    procédure~; effet attendu~; risque résiduel.

  \item \textbf{Mapping ATT\&CK (Analyste Technique)~:}
    À partir de tous les faits disponibles, complétez le tableau
    de mapping MITRE ATT\&CK (minimum 8~entrées)~:
    Timestamp / Événement / Tactique / Technique / TTP ID.

  \item \textbf{Preuves forensiques (Analyste Technique)~:}
    Avant d'éteindre ou de remettre en état, listez les 5~pièces
    à collecter sur srv-erp-01, dans l'ordre de volatilité,
    avec l'outil d'acquisition et le format de stockage.

  \item \textbf{Communication de crise (Responsable Communication)~:}
    Rédigez un email de crise à destination du DG de LiZbiZ-SIS~:
    2~paragraphes, non technique, incluant~: nature de l'incident,
    impact actuel, mesures prises, prochaines étapes,
    décision requise (payer ou ne pas payer).

  \item \textbf{Recommandation de paiement (Incident Manager)~:}
    Faut-il payer la rançon~? Argumentez en 5~points (technique,
    légal, éthique, financier, précédent). Quelle est votre
    recommandation finale~?
\end{enumerate}
\end{boxexo}

\begin{boxlizbiz}
\textbf{Inject~3 — T+1h15~min}

\begin{itemize}
  \item Les sauvegardes de srv-backup-01 ont également été chiffrées~:
    le ransomware a accédé au partage SMB et chiffré les fichiers
    \texttt{.bak} et \texttt{.sql} en priorité.
  \item Bonne nouvelle~: une copie tape hors-site (LTO-7) du dimanche
    13~mai à 23h00 existe dans le local d'archivage (non connecté
    au réseau).
  \item Le partenaire mobile money (MTN MoMo) contacte LiZbiZ par
    email~: il a détecté des transactions suspectes depuis le
    compte API de LiZbiZ entre 04h15 et 04h50.
\end{itemize}
\end{boxlizbiz}

\begin{boxexo}
\textbf{Questions Post-Inject~3~:}
\begin{enumerate}
  \item La sauvegarde tape hors-site est-elle utilisable pour la
    restauration~? Quels risques~? Quelle procédure de validation~?
  \item Les transactions MTN MoMo suspectes modifient-elles les
    obligations légales~? Notifier qui~? Dans quel délai~?
  \item Révisez votre recommandation sur le paiement de la rançon.
\end{enumerate}
\end{boxexo}


% ==============================================================================
\section{Phase~3 (T+2h à T+4h)~: Éradication et rétablissement}
\label{lab3:phase3}
% ==============================================================================

\textbf{Durée allouée~: 30~minutes de simulation.}

\begin{boxexo}
\textbf{Travail Phase~3~:}

\begin{enumerate}
  \item \textbf{Plan d'éradication~:}
    Rédigez le plan d'éradication en 7~étapes ordonnées. Chaque étape~:
    action, responsable, outil, critère de validation, durée estimée.

  \item \textbf{Plan de restauration~:}
    Définissez l'ordre de restauration des systèmes~:
    (a)~quel système restaurer en premier~? Pourquoi~?
    (b)~comment valider l'intégrité de la bande avant restauration~?
    (c)~quelle surveillance renforcer pendant les 30~premiers jours~?

  \item \textbf{Critères de retour à la normale~:}
    Listez 5~critères mesurables qui permettront de déclarer
    la fin de l'incident et le retour à la normale.
\end{enumerate}
\end{boxexo}


% ==============================================================================
\section{Phase~4 — RETEX de l'équipe}
\label{lab3:retex}
% ==============================================================================

À l'issue de la simulation, chaque équipe conduit son propre RETEX~:

\begin{boxexo}
\textbf{RETEX Tabletop — Questions réflexives~:}

\begin{enumerate}
  \item \textbf{Ce qui a bien fonctionné~:}
    Identifiez 3~décisions ou actions que votre équipe a bien gérées.
    Pourquoi~? Quel principe ou outil a été déterminant~?

  \item \textbf{Ce qui aurait pu mieux se passer~:}
    Identifiez 3~décisions retardées ou erronées. Quelle information
    manquait~? Quelle procédure aurait dû exister~?

  \item \textbf{Rôles et communication~:}
    La répartition des rôles était-elle adaptée~? L'Incident Manager
    a-t-il eu trop/pas assez d'informations~? La communication interne
    a-t-elle été fluide~?

  \item \textbf{Ce que LiZbiZ-SIS doit créer avant le prochain incident~:}
    Listez les 5~artefacts documentaires (playbooks, politiques,
    contacts, procédures) dont l'absence vous a pénalisés.

  \item \textbf{Autoévaluation~:}
    Sur une échelle de 1 à~5, évaluez la performance de votre équipe
    sur~: (a)~rapidité de détection~; (b)~qualité du confinement~;
    (c)~communication~; (d)~application des normes~;
    (e)~plan de remédiation. Justifiez chaque note.
\end{enumerate}
\end{boxexo}


% ==============================================================================
\section{Livrable et critères d'évaluation}
\label{lab3:livrable}
% ==============================================================================

\textbf{Rapport d'incident (15~pages) + Présentation orale (10~min).}

Structure du rapport~:
\begin{enumerate}
  \item Résumé exécutif à destination du DG (1~page).
  \item Chronologie de l'incident (tableau annoté MITRE).
  \item Analyse de l'attaque (vecteur, propagation, impact).
  \item Preuves collectées (liste avec hashes et méthode).
  \item Décisions clés et leur justification.
  \item Plan d'éradication et de restauration.
  \item Plan de remédiation priorisé (6~actions minimum, tableau
        impact / effort / délai / responsable / KPI).
  \item RETEX (causes profondes, 5~Pourquoi, leçons).
\end{enumerate}

\begin{table}[H]
\centering\small
\caption{Barème Lab~3}
\label{tab:lab3-bareme}
\rowcolors{2}{TableauPair}{TableauImpair}
\begin{tabular}{p{6cm}p{6.5cm}p{2cm}}
\toprule
\tableauHeader{Critère} & \tableauHeader{Indicateurs} &
\tableauHeader{Points} \\
\midrule
Qualification et premières décisions & Gravité justifiée, actions prioritaires correctes & 1.5 \\
Confinement technique & Actions techniques correctes et ordonnées & 1.5 \\
Mapping MITRE ATT\&CK & $\geq 8$~entrées, TTPs corrects & 1.5 \\
Communication (DG + employés) & Ton adapté, contenu juste, non technique & 1 \\
Traitement légal & Obligations correctes, délais et destinataires & 1 \\
RETEX et plan de remédiation & Causes profondes, 5~Pourquoi, $\geq 6$~actions & 2 \\
Présentation orale & Clarté, maîtrise du sujet, gestion du temps & 1.5 \\
\bottomrule
\multicolumn{2}{r}{\textbf{Total}} & \textbf{10} \\
\end{tabular}
\end{table}
